\section{Hidden Markov Models: A Latent-Variable-Model Definition}
\label{sec:hmm-definition}

Section 6 ("Markov Chains, Hidden Markov Models, and the Markov Blanket")
of our GWP1 report introduced the
Hidden Markov Model (HMM) following Rabiner's classical 5-tuple notation.
Here we give our own definition, grounded instead in the
\emph{latent-variable-model} perspective of Bishop
\autocite{bishop2006prml}, since this framing is what directly motivates
the Expectation-Maximization (Baum-Welch) training algorithm implemented
in Section~\ref{sec:algorithms}.

\subsection{An HMM as a Chain of Latent Variables}

A Hidden Markov Model is a sequence model over a chain of discrete latent
variables $z_1, \ldots, z_T$ (the hidden \emph{regimes}) and a
corresponding chain of observed variables $x_1, \ldots, x_T$ (here, daily
log returns), with the joint distribution factorised as
\[
  p(x_{1:T}, z_{1:T}) = p(z_1)\, p(x_1 \mid z_1) \prod_{t=2}^{T}
    p(z_t \mid z_{t-1})\, p(x_t \mid z_t).
\]
Two conditional-independence assumptions give this factorisation its
name and its tractability: (i) the \textbf{Markov assumption} on the
latent chain, $p(z_t \mid z_{1:t-1}) = p(z_t \mid z_{t-1})$ --- the
regime at time $t$ depends on the past only through the regime at
$t-1$; and (ii) the \textbf{emission independence assumption}, that the
observation $x_t$ depends on the full latent history only through the
current regime $z_t$. Compared to a simple finite mixture model (which
would model $x_{1:T}$ as independent draws from one of $K$ component
distributions), an HMM differs in exactly one respect: the component
("regime") responsible for each observation is itself generated by a
Markov chain rather than drawn independently every period. This single
change is precisely what lets the model capture \textbf{regime
persistence} --- the empirical fact, documented in our GWP1 stylized
facts, that oil-return volatility is strongly autocorrelated (squared
daily-return autocorrelation of 0.46 at lag 1 for WTI) rather than i.i.d.

In our application each latent state $z_t \in \{1, \ldots, K\}$ ($K=3$)
represents a market regime (bear/stagnant/bull), and each emission
distribution $p(x_t \mid z_t = k) = \mathcal{N}(x_t; \mu_k, \sigma_k^2)$
is a univariate Gaussian over that day's log return --- a continuous
emission density, rather than the discrete-symbol alphabet $\Sigma$ used
in the purely combinatorial presentations of the model. This is a
standard and well-documented extension (the "Gaussian HMM"): the
emission distribution is swapped for a parametric density appropriate to
the data type, while the latent-chain structure, and therefore every one
of the four classical algorithms (Section~\ref{sec:algorithms}), is
unchanged in form -- every place a discrete emission probability
$b_j(o_t)$ appears, a Gaussian density $\mathcal{N}(x_t; \mu_j,
\sigma_j^2)$ appears instead.

\subsection{The Complete-Data Likelihood and Why EM Is Needed}

Writing $z_{tk} = \mathbb{1}[z_t = k]$, the complete-data log-likelihood
(if the regimes were observed) would be
\[
  \log p(x_{1:T}, z_{1:T}) = \sum_{k=1}^{K} z_{1k} \log \pi_k
    + \sum_{t=2}^{T} \sum_{j=1}^{K} \sum_{k=1}^{K} z_{t-1,j} z_{tk} \log A_{jk}
    + \sum_{t=1}^{T} \sum_{k=1}^{K} z_{tk} \log \mathcal{N}(x_t; \mu_k, \sigma_k^2),
\]
which would be trivial to maximise in closed form (it decomposes into
independent counting/averaging problems). The difficulty is that the
regimes are never observed: only $x_{1:T}$ is. Baum-Welch is exactly the
Expectation-Maximization (EM) algorithm applied to this latent-variable
model \autocite{bishop2006prml}: the E-step replaces every unobserved
indicator $z_{tk}$ by its posterior expectation $\gamma_t(k) =
\mathbb{E}[z_{tk} \mid x_{1:T}]$ (computed via the Forward-Backward
algorithm, Section~\ref{sec:algorithms}), and the M-step maximises the
\emph{expected} complete-data log-likelihood with those soft
assignments plugged in -- which reduces, again, to simple weighted
counting/averaging. This is precisely why Baum-Welch inherits EM's
general guarantee that the observed-data log-likelihood never decreases
between iterations, which we verify empirically on our toy example in
Section~\ref{sec:algorithms}.
